\chapter{Finiteness descent from the residue field (Stacks 04GG fragment)}
\label{ch:Proetale_Mathlib_RingTheory_Module_FiniteResidue}

This chapter extracts a Mathlib-PR-quality helper for the
``finiteness descends from the residue field'' fragment of
Stacks Tag 04GG. The setting is a Noetherian local ring \(A\) with
maximal ideal \(\mathfrak{m}\) and residue field \(k = A/\mathfrak{m}\),
together with a finitely-presented flat \(A\)-module \(M\) whose
residue \(M \otimes_A k\) is finite-dimensional as a \(k\)-vector
space. The headline theorem is that \(M\) itself is module-finite
over \(A\).

\section{Nakayama in residue form}

\begin{lemma}
  \label{lem:nakayama-residue-spanning}
  \lean{Module.eq_top_of_residue_eq_top_of_finite}
  \leanok
  [Nakayama from a residue spanning set]
  Let \(A\) be a local ring with maximal ideal \(\mathfrak{m}\), and
  let \(M\) be a finitely-generated \(A\)-module. If
  \(M = \mathfrak{m} \cdot M + N\) for some submodule
  \(N \subseteq M\), then \(N = M\).
\end{lemma}

\begin{proof}
  \leanok
  Apply Nakayama's lemma to the quotient \(M/N\). The hypothesis
  \(M = \mathfrak{m} \cdot M + N\) implies
  \(\mathfrak{m} \cdot (M/N) = M/N\): given any \(\bar x \in M/N\)
  with lift \(x \in M\), write \(x = m + n\) with \(m \in
  \mathfrak{m} \cdot M\) and \(n \in N\); then
  \(\bar x = \overline{m} \in \mathfrak{m} \cdot (M/N)\). Since
  \(M\) is finitely generated, so is the quotient \(M/N\). The
  maximal ideal \(\mathfrak{m}\) of a local ring is contained in
  its Jacobson radical, so Nakayama's lemma yields \(M/N = 0\),
  i.e.\ \(N = M\).
\end{proof}

\section{Lifting a residue basis to a partial spanning set}

\begin{lemma}
  \label{lem:lift-residue-basis-to-spanning}
  \lean{Module.span_residue_basis_add_maximal_smul_eq_top}
  \leanok
  [Lift of residue basis to spanning modulo \(\mathfrak{m} M\)]
  Let \(A\) be a local ring with maximal ideal \(\mathfrak{m}\)
  and residue field \(k = A/\mathfrak{m}\), and let \(M\) be an
  \(A\)-module. Let \(b : I \to M\) be a family whose image under
  the projection \(M \twoheadrightarrow M / \mathfrak{m} \cdot M\)
  spans \(M / \mathfrak{m} \cdot M\) as an \(A\)-module. Then
  \[
    \mathrm{span}_A \{b_i\}_{i \in I} \;+\; \mathfrak{m} \cdot M
    \;=\; M.
  \]
\end{lemma}

\begin{proof}
  \leanok
  Fix \(x \in M\) and consider its image \(\bar x \in M / \mathfrak{m} \cdot M\).
  By hypothesis, \(\bar x\) lies in the image of
  \(\mathrm{span}_A \{b_i\}_{i \in I}\) under the projection,
  so there exists \(y \in \mathrm{span}_A \{b_i\}_{i \in I}\) with
  \(\bar y = \bar x\). Then \(x - y \in \mathfrak{m} \cdot M\), and
  \(x = y + (x - y) \in \mathrm{span}_A \{b_i\}_{i \in I}
  + \mathfrak{m} \cdot M\).
\end{proof}

\section{Headline: finiteness descent}

\begin{theorem}
  \label{thm:finite-of-finite-residue-flat-local}
  \lean{Module.finite_of_finite_residue_of_flat_of_local}
  \leanok
  [Stacks 04GG, finiteness descent fragment]
  Let \(A\) be a local ring with maximal ideal \(\mathfrak{m}\),
  and let \(M\) be a finitely-presented flat \(A\)-module. If
  \(M / \mathfrak{m} \cdot M\) is finite as an \(A\)-module, then
  \(M\) is module-finite over \(A\).
\end{theorem}

\begin{proof}
  \leanok
  \uses{lem:lift-residue-basis-to-spanning, lem:nakayama-residue-spanning}
  \emph{Step 1: choose a residue generating set and a lift.}
  Pick a finite generating set \(\bar s\) of
  \(M / \mathfrak{m} \cdot M\), and choose for each element of
  \(\bar s\) a lift in \(M\); call \(t \subseteq M\) the resulting
  finite set.

  \emph{Step 2: residue surjectivity.} By
  \cref{lem:lift-residue-basis-to-spanning},
  \(\mathrm{span}_A t + \mathfrak{m} \cdot M = M\).

  \emph{Step 3: apply Nakayama.} By
  \cref{lem:nakayama-residue-spanning} applied to \(M\) and the
  submodule \(N := \mathrm{span}_A t\), the hypothesis
  \(M = \mathfrak{m} \cdot M + N\) from Step~2 yields \(N = M\).

  \emph{Conclusion.} The finite set \(t\) generates \(M\) as an
  \(A\)-module. This is exactly \(\mathrm{Module.Finite}\,A\,M\).

  \emph{Remark on the flatness hypothesis.} The argument above
  uses only that \(M\) is finitely-presented (which entails
  finite-generation of \(M\)) and Nakayama. Flatness of \(M\) is
  the natural hypothesis in Stacks~04GG and is retained in the
  headline statement for direct alignment with downstream
  applications (where \(M\) is the underlying module of an étale
  \(A\)-algebra, hence automatically flat); the proof itself does
  not exploit flatness.
\end{proof}

\section{Étale specialisation}

\begin{corollary}
  \label{cor:finite-of-finite-residue-etale-local}
  \lean{Algebra.Etale.finite_of_finite_residue_of_local}
  \leanok
  [Stacks 04GG, étale fibre form]
  Let \(A\) be a local ring with maximal ideal \(\mathfrak{m}\),
  and let \(B\) be an étale \(A\)-algebra whose underlying
  \(A\)-module is finitely-presented and such that
  \(B / \mathfrak{m} \cdot B\) is finite as an \(A\)-module. Then
  \(B\) is module-finite over \(A\).
\end{corollary}

\begin{proof}
  \leanok
  \uses{thm:finite-of-finite-residue-flat-local}
  An étale \(A\)-algebra is flat (via
  \texttt{Algebra.Etale.flat}). Applying
  \cref{thm:finite-of-finite-residue-flat-local} to the underlying
  \(A\)-module \(M := B\) yields \(\mathrm{Module.Finite}\,A\,B\).
\end{proof}

\section{References}

\begin{itemize}
\item Stacks Project, Tag~\href{https://stacks.math.columbia.edu/tag/04GG}{04GG} —
  the at-the-fibre / finiteness descent statement: a finitely
  presented flat module over a Noetherian local ring whose special
  fibre is finite-dimensional is module-finite. The primary source
  of \cref{thm:finite-of-finite-residue-flat-local} and
  \cref{cor:finite-of-finite-residue-etale-local}.
\item Stacks Project, Tag~\href{https://stacks.math.columbia.edu/tag/00DV}{00DV} —
  Nakayama's lemma in module form. The source of
  \cref{lem:nakayama-residue-spanning} (in the local-ring
  specialisation).
\end{itemize}
